\setlength{\parindent}{2em}

\begin{center}
\textbf{\large{Tóm tắt}	}
\end{center}

\addcontentsline{toc}{chapter}{Tóm tắt}

\begin{small}

Trong bối cảnh chuyển đổi số mạnh mẽ, các cuộc tấn công lừa đảo (phishing) vẫn là một trong những rủi ro an ninh mạng phổ biến. Kẻ tấn công liên tục thay đổi chiến lược sử dụng các kỹ thuật che giấu như rút gọn link (URL shortening), tạo tên miền nhái (Typosquatting) khiến các hệ thống phòng thủ tĩnh trở nên kém hiệu quả. Cần một giải pháp bảo vệ trực tiếp tại phía người dùng (Client-side), có khả năng phân tích và đưa ra cảnh báo ngay lập tức đối với các URL đáng ngờ. Dự án CheckPost được phát triển nhằm giải quyết vấn đề bằng cách xây dựng tiện ích trình duyệt (Manifest V3) tích hợp trí tuệ nhân tạo. Đặc biệt, nghiên cứu này đã khắc phục triệt để hiện tượng rò rỉ dữ liệu (Data Leakage) thường gặp trong học máy bảo mật, đồng thời tối ưu hóa thuật toán XGBoost kết hợp cơ chế kiểm định Heuristic (Khoảng cách Levenshtein) để phân tích đa tầng các đặc trưng URL. Phương pháp này cho phép hệ thống đưa ra dự đoán với tỷ lệ chặn nhầm (False Positive Rate) tiệm cận 0\%, cùng với khả năng giải thích AI (Explainable AI bằng SHAP), hỗ trợ ngăn chặn hiệu quả các cuộc tấn công Zero-day. Kết quả thực nghiệm khẳng định tiện ích hoạt động ổn định với độ trễ cực thấp, đảm bảo an toàn thông tin mà không gây ảnh hưởng đến hiệu năng duyệt web.

\vspace*{1cm}
\textbf{Từ khóa}: An ninh mạng, Tiện Ích Trình Duyệt, Phishing, XGBoost, Explainable AI, Data Leakage, Manifest V3

    

\end{small}